\documentclass[11pt]{article}

% ============================================================
% Packages
% ============================================================
\usepackage[utf8]{inputenc}
\usepackage[margin=1in]{geometry}
\usepackage{amsmath,amssymb,amsthm}
\usepackage{graphicx}
\usepackage{booktabs}
\usepackage{xcolor}
\usepackage{hyperref}
\usepackage{tikz}
\usepackage{enumitem}
\usepackage{pgfplots}

\pgfplotsset{compat=1.18}

% ============================================================
% Document Formatting
% ============================================================
\setlength{\parskip}{0.6em}
\setlength{\parindent}{0pt}

% ============================================================
% List Formatting
% ============================================================
\setlist[enumerate]{
    leftmargin=2em,
    itemsep=0.25em
}

% ============================================================
% Solution Environment
% ============================================================
\newenvironment{solution}
    {\par\noindent\textbf{Solution.}\par}
    {\par}

% ============================================================
% Link Styling
% ============================================================
\hypersetup{
    colorlinks=true,
    linkcolor=blue,
    urlcolor=cyan
}

% ============================================================
% Metadata
% ============================================================
\title{Ordinary Differential Equations: Abridged Notes}
\author{Chris Ritter}
\date{October 13, 2026}

% ============================================================
% Document
% ============================================================

\begin{document}
\maketitle
\begin{center}
\small\emph{These abridged notes were reconstructed afterward because I did not take
direct notes during class. They are not a record of the actual lecture.}
\end{center}
\tableofcontents
\newpage

\section{Topic Summary}
A repeated characteristic root supplies fewer distinct exponential
solutions than the order of the equation requires. For a double real
root \(r\), the second independent solution is \(te^{rt}\).
Reduction of order is a more general method: when one nonzero
solution \(y_1\) of a second-order homogeneous linear equation is
known, seek another as \(y_2=v(t)y_1(t)\). This method also works
when coefficients vary with time.

\section{Repeated Roots}
If \(ar^2+br+c=a(r-r_0)^2\), then
\[
\boxed{y(t)=(C_1+C_2t)e^{r_0t}.}
\]
The factor \(t\) makes \(te^{r_0t}\) independent of
\(e^{r_0t}\). The two are a fundamental pair because their
Wronskian is \(e^{2r_0t}\ne0\).

\subsection{Example: A Repeated-Root Initial-Value Problem}
\textbf{Adapted from Boyce et al., 11th ed., \S3.4, Problem 9.}
Solve
\[
9y''-12y'+4y=0,\qquad y(0)=2,\qquad y'(0)=-1.
\]
\begin{solution}
The characteristic polynomial is \((3r-2)^2\), so the double root
is \(r_0=2/3\). Write \(y=(C_1+C_2t)e^{2t/3}\). The first initial
condition gives \(C_1=2\). Since
\[
y'=e^{2t/3}\left(C_2+\frac23(C_1+C_2t)\right),
\]
the second gives \(C_2+4/3=-1\), hence \(C_2=-7/3\). Therefore
\[
\boxed{y(t)=\left(2-\frac73t\right)e^{2t/3}.}
\]
\end{solution}

\section{Reduction of Order}
For a normalized equation \(y''+p(t)y'+q(t)y=0\) and a known
solution \(y_1\) that is nonzero on the interval of work, substitute
\(y_2=vy_1\). Simplifying with the equation satisfied by \(y_1\)
yields a first-order equation for \(v'\). One convenient formula is
\[
\boxed{y_2(t)=y_1(t)\int
\frac{\exp\left(-\int p(t)\,dt\right)}{y_1(t)^2}\,dt.}
\]
Any constant added inside the outer integral merely adds a multiple
of \(y_1\); choose the constant to obtain an independent solution.

\subsection{Example: Variable Coefficients}
\textbf{Adapted from Boyce et al., 11th ed., \S3.4, Problem 21.}
Given \(y_1(x)=\sin(x^2)\), find a second solution of
\[
xy''-y'+4x^3y=0,\qquad x>0.
\]
\begin{solution}
After division by \(x\), \(p(x)=-1/x\), so
\(e^{-\int p\,dx}=x\) up to a constant factor. On an interval
where \(\sin(x^2)\ne0\), reduction of order gives
\[
\begin{aligned}
y_2
&=\sin(x^2)\int\frac{x}{\sin^2(x^2)}\,dx\\
&=\frac12\sin(x^2)\int\csc^2u\,du,
\qquad u=x^2\\
&=-\frac12\cos(x^2).
\end{aligned}
\]
A nonzero constant multiple is equally valid, so choose
\[
\boxed{y_2(x)=\cos(x^2).}
\]
This expression solves the equation for every \(x>0\), including
zeros of \(\sin(x^2)\); the intermediate integral is used locally
where its denominator is nonzero. The pair is independent because
its Wronskian is \(-2x\ne0\) for \(x>0\).
\end{solution}
\end{document}
